\documentclass[12pt, a4paper]{article}

% Packages
\usepackage[margin=1in]{geometry}
\usepackage{graphicx}
\usepackage{tabularx}
\usepackage{hyperref}
\usepackage{enumitem}
\usepackage{booktabs}
\usepackage{float}

% Title and Author
\title{\textbf{Gen AI Driven Wellbeing Coach and Health Management Platform} \\ \large Software Requirements Specification}
\author{\textbf{Team 48} \\ Rishabh Goyal, Aditya Singh, Tanush Garg, Avnish Uba, Sashi Kumar Gudipati}
\date{}

\begin{document}
\raggedbottom

\maketitle

\section{Brief Problem Statement}
Health data is fragmented across wearable apps and medical reports (often PDFs), making it hard for individuals to understand their health status and take timely action. Clinicians also lack a lightweight way to monitor patients between visits without wading through raw, unstructured data. This project builds a platform that ingests medical reports, extracts verifiable facts, and provides explainable alerts and AI-assisted guidance grounded in the user’s own evidence.

\section{System Requirements (Technologies/Languages/Libraries)}
\begin{itemize}[leftmargin=*]
    \item \textbf{Backend:} Python 3, FastAPI, Uvicorn
    \item \textbf{Storage + DB:} Supabase (PostgreSQL + Storage), Supabase Python client
    \item \textbf{OCR + PDF processing:} Tesseract OCR (\texttt{pytesseract}), \texttt{pdf2image}, Pillow, OpenCV, NumPy
    \item \textbf{OS dependencies:} Tesseract installed and Poppler available for PDF-to-image conversion
    \item \textbf{Text preprocessing:} OCR text cleaning + chunking utilities; LangChain text splitters (for recursive chunking)
    \item \textbf{Embeddings (RAG):} \texttt{SentenceTransformers} with model \texttt{BAAI/bge-base-en-v1.5}
    \item \textbf{Configuration:} Environment variables + \texttt{.env} support (e.g., \texttt{SUPABASE\_URL}, \texttt{SUPABASE\_SERVICE\_ROLE\_KEY}, bucket/table names)
    \item \textbf{(Planned for R2 / integration):} External context APIs (weather/AQI), LLM reasoning layer (Gemini) for coached answers with citations
\end{itemize}

\section{Users Profile}
\begin{itemize}[leftmargin=*]
    \item \textbf{Patient (primary):} Uploads medical reports; views extracted labs, alerts, and AI coach explanations. Comfortable using a web/mobile UI, not necessarily medically trained.
    \item \textbf{Doctor/Clinician (secondary):} Reviews a list of patients and risk status; drills into evidence-backed alerts and summaries. High computer proficiency; expects concise, explainable outputs.
    \item \textbf{System Admin/Developer (support):} Configures Supabase, environment variables, and monitors ingestion/OCR failures. Technical user.
\end{itemize}

\section{Requirements}

\subsection{Feature Requirements (Use Cases)}

\begin{table}[htbp]
    \centering
    \small
    \renewcommand{\arraystretch}{1.2}
    \begin{tabularx}{\textwidth}{|>{\raggedright\arraybackslash}p{4cm}|X|c|}
        \hline
        \textbf{Feature / Use Case} & \textbf{Description} & \textbf{Release} \\
        \hline
        Upload medical report & Patient uploads a PDF/image report which is stored under the patient’s scope in secure storage; system returns a reference path/URL. & R1 \\
        \hline
        View uploaded report metadata & Patient can view basic metadata for an uploaded report (filename, upload timestamp, link/reference). & R1 \\
        \hline
        Run OCR on a stored report & Patient triggers OCR; system downloads the stored file, performs OCR, and persists OCR text + confidence. & R1 \\
        \hline
        View OCR output & Patient can view OCR text and an OCR confidence/quality indicator for the report. & R1 \\
        \hline
        Extract structured lab results & Patient triggers deterministic (regex-based) extraction of lab measurements from OCR text and stores them as atomic facts linked to the report. & R1 \\
        \hline
        View structured lab results & Patient can list lab results for a report including value, unit, reference range, and source page when available. & R1 \\
        \hline
        Clean OCR text for retrieval & System cleans noisy OCR text to remove repeated headers/junk and produce meaningful lines for downstream retrieval. & R1 \\
        \hline
        Chunk OCR text & System splits cleaned text into retrieval-ready chunks (table-like rows + narrative chunks) to support evidence retrieval. & R1 \\
        \hline
        Embed chunks & System generates vector embeddings for chunks using a configured embedding model for semantic search. & R1 \\
        \hline
        Retrieve evidence for a question & Patient asks a question; system retrieves top relevant chunks with source metadata (file + page) as citations. & R1 \\
        \hline
        AI coach answer with citations & Patient receives an answer that is grounded in retrieved evidence and returns citations/snippets for transparency. & R1 \\
        \hline
        Generate explainable alerts & System produces deterministic alerts (low/medium/high) from extracted facts and links each alert to evidence (report/lab/snippet). & R1 \\
        \hline
        View dashboard summary & Patient views a dashboard with wellbeing score/trend (and optionally environment context) plus active alert counts. & R1 \\
        \hline
        Doctor dashboard: list patients & Doctor views their patient list with risk level and last updated timestamps. & R2 \\
        \hline
    \end{tabularx}
\end{table}

\subsection{Use Case Diagram}
\textit{Actors to include: Patient, Doctor, System Admin/Developer. External systems: Supabase Storage/DB, OCR engine (Tesseract), Weather/AQI API, LLM provider.}

\begin{figure}[htbp]
    \centering
    \includegraphics[width=\textwidth, keepaspectratio]{image.jpeg}
    \caption{System Use Case Diagram}
    \label{fig:usecase}
\end{figure}

\section{Use Case Descriptions}

\subsection*{UC-01: Upload medical report}
\textbf{Overview:} Patient uploads a medical report (PDF/image). The system stores the file in secure storage under the patient scope and returns a reference. \\
\textbf{Actors:} Patient \\
\textbf{Pre-condition:} Patient has a valid user identifier; system storage is configured. \\
\textbf{Flow:}
\begin{itemize}[leftmargin=*]
    \item \textbf{Main (success) Flow:}
    \begin{enumerate}
        \item Patient selects a report file and submits upload.
        \item System validates non-empty file and required user identifier.
        \item System generates a unique, user-scoped storage path.
        \item System uploads bytes to storage with the correct content type.
        \item System returns storage reference (path) and accessible URL (public/signed per deployment).
    \end{enumerate}
    \item \textbf{Alternate Flows:}
    \begin{itemize}
        \item \textit{A1.} Missing user identifier $\rightarrow$ System rejects request with validation error. Post-condition: no file stored.
        \item \textit{A2.} Empty/corrupt file $\rightarrow$ System rejects request. Post-condition: no file stored.
        \item \textit{A3.} Storage failure $\rightarrow$ System returns error. Post-condition: no file stored.
    \end{itemize}
\end{itemize}
\textbf{Post-Condition:} Report file is stored and uniquely addressable; patient receives reference.

\subsection*{UC-02: View uploaded report metadata}
\textbf{Overview:} Patient views metadata for an uploaded report to confirm it is available for processing. \\
\textbf{Actors:} Patient \\
\textbf{Pre-condition:} Patient has uploaded at least one report. \\
\textbf{Flow:}
\begin{itemize}[leftmargin=*]
    \item \textbf{Main Flow:}
    \begin{enumerate}
        \item Patient opens “My Reports”.
        \item System lists reports with filename, upload time, and processing status (uploaded/ocr\_done/extracted).
    \end{enumerate}
    \item \textbf{Alternate Flows:}
    \begin{itemize}
        \item \textit{A1.} No reports $\rightarrow$ System shows empty state.
    \end{itemize}
\end{itemize}
\textbf{Post-Condition:} Patient can identify which report to OCR/extract.

\subsection*{UC-03: Run OCR on a stored report}
\textbf{Overview:} Patient triggers OCR on a previously uploaded report. System performs OCR and persists OCR text and confidence. \\
\textbf{Actors:} Patient, System (OCR engine) \\
\textbf{Pre-condition:} Report exists in storage; OCR dependencies available. \\
\textbf{Flow:}
\begin{itemize}[leftmargin=*]
    \item \textbf{Main Flow:}
    \begin{enumerate}
        \item Patient selects a report and starts OCR.
        \item System downloads the file bytes from storage.
        \item If PDF, system converts pages to images.
        \item System preprocesses images for OCR quality.
        \item System runs OCR and computes average confidence.
        \item System stores OCR text + confidence and source metadata in the database.
    \end{enumerate}
    \item \textbf{Alternate Flows:}
    \begin{itemize}
        \item \textit{A1.} Storage path invalid/not found $\rightarrow$ System returns “report not found”.
        \item \textit{A2.} PDF conversion fails (missing Poppler, corrupt PDF) $\rightarrow$ System returns OCR error; status marked failed.
        \item \textit{A3.} OCR engine fails $\rightarrow$ System returns OCR error; status marked failed.
    \end{itemize}
\end{itemize}
\textbf{Post-Condition:} OCR output and confidence are stored for the report (or failure is recorded).

\subsection*{UC-04: View OCR output}
\textbf{Overview:} Patient reviews OCR text and confidence to validate scan quality before extraction/AI queries. \\
\textbf{Actors:} Patient \\
\textbf{Pre-condition:} OCR has been completed for the report. \\
\textbf{Flow:}
\begin{itemize}[leftmargin=*]
    \item \textbf{Main Flow:}
    \begin{enumerate}
        \item Patient opens report OCR results.
        \item System displays OCR text and confidence indicator.
    \end{enumerate}
    \item \textbf{Alternate Flows:}
    \begin{itemize}
        \item \textit{A1.} OCR not available $\rightarrow$ System asks patient to run OCR first.
    \end{itemize}
\end{itemize}
\textbf{Post-Condition:} Patient can confirm OCR readiness.

\subsection*{UC-05: Extract structured lab results}
\textbf{Overview:} Patient triggers deterministic extraction of lab values from OCR text and stores each lab as an atomic fact linked to the report. \\
\textbf{Actors:} Patient, System \\
\textbf{Pre-condition:} OCR text exists for the report. \\
\textbf{Flow:}
\begin{itemize}[leftmargin=*]
    \item \textbf{Main Flow:}
    \begin{enumerate}
        \item Patient clicks “Extract Labs”.
        \item System loads OCR text for the report.
        \item System applies allow-listed test name matching and numeric/unit parsing.
        \item System stores extracted lab results linked to the report.
    \end{enumerate}
    \item \textbf{Alternate Flows:}
    \begin{itemize}
        \item \textit{A1.} Report not found $\rightarrow$ System returns error; no rows inserted.
        \item \textit{A2.} No extractable labs $\rightarrow$ System returns success with 0 inserted.
        \item \textit{A3.} DB insert fails $\rightarrow$ System returns error; partial inserts handled per transaction strategy.
    \end{itemize}
\end{itemize}
\textbf{Post-Condition:} Lab results are stored (or the system reports that none were confidently extractable).

\subsection*{UC-06: View structured lab results}
\textbf{Overview:} Patient views extracted labs with units, reference ranges, and traceability (source page/snippet when available). \\
\textbf{Actors:} Patient \\
\textbf{Pre-condition:} Labs have been extracted for the report. \\
\textbf{Flow:}
\begin{itemize}[leftmargin=*]
    \item \textbf{Main Flow:}
    \begin{enumerate}
        \item Patient opens the “Labs” view for a report.
        \item System lists lab results with test name, numeric value, unit, reference range, and provenance.
    \end{enumerate}
    \item \textbf{Alternate Flows:}
    \begin{itemize}
        \item \textit{A1.} No labs extracted yet $\rightarrow$ System prompts to run extraction.
    \end{itemize}
\end{itemize}
\textbf{Post-Condition:} Patient can inspect verifiable facts extracted from the report.

\subsection*{UC-07: Clean OCR text for retrieval}
\textbf{Overview:} System transforms raw OCR into cleaned lines by removing repeated headers/junk and preserving medically relevant content. \\
\textbf{Actors:} System \\
\textbf{Pre-condition:} OCR text exists. \\
\textbf{Flow:}
\begin{itemize}[leftmargin=*]
    \item \textbf{Main Flow:}
    \begin{enumerate}
        \item System splits raw text into line-like segments.
        \item System removes known junk patterns (headers, page markers, boilerplate).
        \item System returns cleaned text suitable for chunking/retrieval.
    \end{enumerate}
    \item \textbf{Alternate Flows:}
    \begin{itemize}
        \item \textit{A1.} OCR text empty $\rightarrow$ System returns empty cleaned result.
    \end{itemize}
\end{itemize}
\textbf{Post-Condition:} Cleaned text is produced without altering the stored immutable OCR original.

\subsection*{UC-08: Chunk OCR text}
\textbf{Overview:} System produces retrieval-ready chunks (tabular rows + narrative chunks) for semantic search and citation. \\
\textbf{Actors:} System \\
\textbf{Pre-condition:} Cleaned OCR text exists. \\
\textbf{Flow:}
\begin{itemize}[leftmargin=*]
    \item \textbf{Main Flow:}
    \begin{enumerate}
        \item System detects section-like headings and splits into sections.
        \item System extracts measurement-like rows as chunks.
        \item System recursively splits remaining narrative text into bounded chunks with overlap.
    \end{enumerate}
\end{itemize}
\textbf{Post-Condition:} A chunk set exists for embedding/retrieval.

\subsection*{UC-09: Embed chunks}
\textbf{Overview:} System generates vector embeddings for chunks for semantic retrieval. \\
\textbf{Actors:} System \\
\textbf{Pre-condition:} Chunk set exists; embedding model available. \\
\textbf{Flow:}
\begin{itemize}[leftmargin=*]
    \item \textbf{Main Flow:}
    \begin{enumerate}
        \item System loads the configured embedding model.
        \item System embeds chunks and stores vectors + required metadata (user/report/source).
    \end{enumerate}
    \item \textbf{Alternate Flows:}
    \begin{itemize}
        \item \textit{A1.} Embedding model not available $\rightarrow$ System fails gracefully and logs error.
    \end{itemize}
\end{itemize}
\textbf{Post-Condition:} Chunks are searchable by vector similarity.

\subsection*{UC-10: Retrieve evidence for a question}
\textbf{Overview:} Patient asks a health question; system retrieves the most relevant chunks with source metadata for citations. \\
\textbf{Actors:} Patient, System \\
\textbf{Pre-condition:} Chunks are embedded for the patient’s documents (or a fallback retrieval mode exists). \\
\textbf{Flow:}
\begin{itemize}[leftmargin=*]
    \item \textbf{Main Flow:}
    \begin{enumerate}
        \item Patient submits a question.
        \item System embeds the query.
        \item System retrieves top-k relevant chunks limited to the patient scope.
        \item System returns retrieved chunks with source file and page metadata.
    \end{enumerate}
    \item \textbf{Alternate Flows:}
    \begin{itemize}
        \item \textit{A1.} No vectors available $\rightarrow$ System returns empty citations and suggests uploading/processing reports.
    \end{itemize}
\end{itemize}
\textbf{Post-Condition:} Evidence set is available for answer generation.

\subsection*{UC-11: AI coach answer with citations}
\textbf{Overview:} System generates an answer that is grounded in retrieved evidence and returns citations/snippets. \\
\textbf{Actors:} Patient, System (LLM) \\
\textbf{Pre-condition:} Evidence retrieved (UC-10). \\
\textbf{Flow:}
\begin{itemize}[leftmargin=*]
    \item \textbf{Main Flow:}
    \begin{enumerate}
        \item System builds a prompt containing the question + retrieved evidence.
        \item LLM generates an answer constrained to evidence.
        \item System returns answer plus citations (source file, page, snippet).
    \end{enumerate}
    \item \textbf{Alternate Flows:}
    \begin{itemize}
        \item \textit{A1.} LLM unavailable $\rightarrow$ System returns error and still shows retrieved evidence.
        \item \textit{A2.} Evidence empty $\rightarrow$ System returns a safe response asking for more data and avoids unsupported claims.
    \end{itemize}
\end{itemize}
\textbf{Post-Condition:} Patient receives an answer and can trace it to sources.

\subsection*{UC-12: Generate explainable alerts}
\textbf{Overview:} System generates deterministic alerts from extracted facts and stores explicit evidence links for every alert. \\
\textbf{Actors:} System \\
\textbf{Pre-condition:} Labs and/or other signals exist. \\
\textbf{Flow:}
\begin{itemize}[leftmargin=*]
    \item \textbf{Main Flow:}
    \begin{enumerate}
        \item System evaluates rules over lab results and trends.
        \item System assigns severity (low/medium/high).
        \item System stores alert plus evidence references (report/lab/snippet).
    \end{enumerate}
    \item \textbf{Alternate Flows:}
    \begin{itemize}
        \item \textit{A1.} Insufficient evidence $\rightarrow$ System does not create alert.
    \end{itemize}
\end{itemize}
\textbf{Post-Condition:} Alerts exist only when explainable via stored evidence.

\subsection*{UC-13: View dashboard summary}
\textbf{Overview:} Patient views a summary of wellbeing score/trend, environment context (if enabled), and active alert counts. \\
\textbf{Actors:} Patient \\
\textbf{Pre-condition:} Patient has data (reports/labs/alerts). \\
\textbf{Flow:}
\begin{itemize}[leftmargin=*]
    \item \textbf{Main Flow:}
    \begin{enumerate}
        \item Patient opens dashboard.
        \item System loads computed summary metrics and active alerts count.
        \item System displays environment context if available.
    \end{enumerate}
    \item \textbf{Alternate Flows:}
    \begin{itemize}
        \item \textit{A1.} New user/no data $\rightarrow$ Dashboard shows onboarding guidance and zero-state metrics.
    \end{itemize}
\end{itemize}
\textbf{Post-Condition:} Patient understands current status at a glance.

\subsection*{UC-14: Doctor dashboard — list patients}
\textbf{Overview:} Doctor views a list of assigned patients with risk level and last update time for prioritization. \\
\textbf{Actors:} Doctor \\
\textbf{Pre-condition:} Doctor is authorized; patients are linked to the doctor. \\
\textbf{Flow:}
\begin{itemize}[leftmargin=*]
    \item \textbf{Main Flow:}
    \begin{enumerate}
        \item Doctor opens patient list.
        \item System returns patients with risk level (low/medium/high) and last-updated timestamp.
    \end{enumerate}
    \item \textbf{Alternate Flows:}
    \begin{itemize}
        \item \textit{A1.} Unauthorized doctor $\rightarrow$ Access denied.
        \item \textit{A2.} No patients assigned $\rightarrow$ Empty list.
    \end{itemize}
\end{itemize}
\textbf{Post-Condition:} Doctor can prioritize reviews.

\end{document}